\chapter{Diseño de la Persistencia de la Información}

\section{Modelo Relacional}
Presentar el modelo entidad-relación o el esquema relacional del sistema.

\begin{figure}[h!]
	\centering
	\fbox{\parbox{0.85\linewidth}{\centering\vspace{3cm}[Diagrama del modelo relacional]\vspace{3cm}}}
	\caption{Modelo relacional de la base de datos.}
\end{figure}

\section{Diccionario de Datos}

\begin{table}[h!]
	\centering
	\caption{Diccionario de datos --- Tabla \texttt{usuarios}}
	\begin{tabular}{>{\ttfamily}l l l l p{5cm}}
		\toprule
		\textbf{Campo} & \textbf{Tipo} & \textbf{Longitud} & \textbf{Nulo} & \textbf{Descripción} \\
		\midrule
		id\_usuario  & INT       & 11  & NO  & Identificador único del usuario (PK). \\
		nombre       & VARCHAR   & 100 & NO  & Nombre completo del usuario. \\
		email        & VARCHAR   & 150 & NO  & Correo electrónico (único). \\
		contrasena   & VARCHAR   & 255 & NO  & Hash de la contraseña. \\
		fecha\_alta  & DATETIME  & --  & NO  & Fecha y hora de registro. \\
		activo       & TINYINT   & 1   & NO  & 1 = activo, 0 = inactivo. \\
		\bottomrule
	\end{tabular}
\end{table}

% Repetir la tabla para cada entidad del modelo

\section{Diseño de las Consultas}

\subsection*{Diccionario de Consultas}

\begin{table}[h!]
	\centering
	\caption{Diccionario de consultas}
	\begin{tabular}{l p{4cm} p{5.5cm}}
		\toprule
		\textbf{ID} & \textbf{Descripción} & \textbf{Sentencia SQL} \\
		\midrule
		Q-001 & Autenticar usuario por email y contraseña &
		\texttt{SELECT * FROM usuarios WHERE email=? AND contrasena=? AND activo=1} \\
		\addlinespace
		Q-002 & Obtener perfil del usuario por ID &
		\texttt{SELECT id\_usuario, nombre, email FROM usuarios WHERE id\_usuario=?} \\
		\addlinespace
		Q-003 & [Descripción de la consulta] & \texttt{[sentencia SQL]} \\
		\bottomrule
	\end{tabular}
\end{table}